\chapter{MCMC Family III: Slice and Elliptical}
\label{ch:mcmc-3}

\newthought{Slice samplers avoid the accept/reject tuning} that makes random-walk \textsc{mcmc}
fiddly. Instead of proposing a step and maybe rejecting it, they sample \emph{under the density
curve} directly, adapting their step size automatically. Both methods here use the standard
\textsc{mcmc} \yamlkey{Bounds} from Chapter~\ref{ch:mcmc-1}.

\section{SliceMCMC}
\label{sec:s-slicemcmc}

\paragraph{How it works} To take a step, slice sampling first draws a height uniformly between zero
and the density at the current point---defining a horizontal ``slice'' through the distribution. It
then samples the next point uniformly from the part of the slice where the density is at least that
height. \Jarvis\ finds that region by stepping outward to bracket an interval (initial width set by
\yamlkey{proposal\_scale}) and then shrinking it until a point inside the slice is found. The step
size therefore adapts to the local shape on its own, with no acceptance rate to tune.

\paragraph{When to use} A robust, low-tuning alternative to random-walk \textsc{mcmc}, especially
when you do not want to hand-fit a proposal scale. \yamlkey{proposal\_scale} only sets the initial
bracket width; the method adjusts from there.

\begin{yamlwin}
Sampling:
  Method: "SliceMCMC"
  Variables: [ ... ]
  Bounds:
    num_chains: 8
    num_iters: 1000
    proposal_scale: 0.5
  LogLikelihood:
    - {name: "LogL", expression: "..."}
\end{yamlwin}

\section{ESS (elliptical slice sampling)}
\label{sec:s-ess}

\paragraph{How it works} Elliptical slice sampling is built for posteriors with a \emph{Gaussian
prior}. From the current point it draws an auxiliary sample from the prior and forms an
\emph{ellipse} through the two; it then slice-samples an angle along that ellipse to pick the next
point. Every point on the ellipse respects the Gaussian prior automatically, so there is no step
size to tune and no rejections to waste---moves are always accepted somewhere on the ellipse.

\paragraph{When to use} Models where the prior is (close to) Gaussian and you want efficient,
tuning-free sampling. It is less suited to flat-prior problems, where the slice methods above or
the Chapter~\ref{ch:mcmc-1} samplers are a better fit.

\begin{yamlwin}
Sampling:
  Method: "ESS"
  Variables: [ ... ]
  Bounds:
    num_chains: 8
    num_iters: 1000
    proposal_scale: 1.0
  LogLikelihood:
    - {name: "LogL", expression: "..."}
\end{yamlwin}

\section{Summary}
\label{sec:mcmc3-summary}

Both samplers trade explicit proposal tuning for automatic step adaptation: \sampler{SliceMCMC} for
general use, \sampler{ESS} when a Gaussian prior makes the elliptical move natural.
